\documentclass[aspectratio=169]{beamer}
\usepackage[portuguese,english]{babel}
\usepackage{slide-tcc}

\usepackage{amsfonts}
\usepackage{xcolor}
\usepackage{datetime}
\usepackage{xmpmulti}

\usepackage[utf8]{inputenc}
\usepackage[T1]{fontenc}
\usepackage{graphicx}
\usepackage{multicol}
\usepackage{subfigure}
\usepackage{lmodern}
\usepackage{amsmath}
\usepackage{mathtools}
\usepackage{mathrsfs}
\usepackage{amssymb}
\usepackage{textcomp}
\usepackage{gensymb}
\usepackage[absolute,overlay]{textpos}
\usepackage{graphicx}
\usepackage{hyperref}
\usepackage{graphicx,url}
\usepackage{subfigure}
\usepackage{latexsym,amsmath,xcolor,multicol,booktabs,calligra}
\usepackage{graphicx,pstricks,listings,stackengine}
\setbeamertemplate{itemize subitem}{$\rightarrow$}
%\usepackage{lipsum}
\usepackage[subnum]{cases}
\usepackage{textpos}
\usepackage{fancybox,calc}
\usepackage{palatino,mathpazo}
\usepackage{amsfonts}
\usepackage{sidecap}
\usepackage{booktabs}
\usepackage{changepage}
\usepackage{wrapfig}
\usepackage{pgfplots}
\usepackage{tikz}
    \usetikzlibrary{shapes.geometric}
\usepackage{xcolor}
\usepackage{datetime2}

\pgfplotsset{compat=1.18, width=7cm}

\def\cmd#1{\texttt{\color{blue}\footnotesize $\backslash$#1}}
\def\env#1{\texttt{\color{cyan}\footnotesize #1}}
\definecolor{deepblue}{rgb}{0,0,0.5}
\definecolor{deepred}{rgb}{0.6,0,0}
\definecolor{deepgreen}{rgb}{0,0.5,0}
\definecolor{halfgray}{gray}{0.55}
\lstset{
    basicstyle=\ttfamily\small,
    keywordstyle=\bfseries\color{deepblue},
    emphstyle=\ttfamily\color{deepred},    % Custom highlighting style
    stringstyle=\color{deepgreen},
    numbers=left,
    numberstyle=\small\color{halfgray},
    rulesepcolor=\color{red!20!green!20!blue!20},
    frame=shadowbox,
}

\title{Trabalho de conclusão de curso}
\subtitle{Utilização de Chain-of-Thought Prompting para classificação de questões da OBI
}
\author{Davi Queiroz - Rodrigo Seiti \\ Profº Amaury Antônio de Castro Junior}
\institute{FACOM \\
UFMS \\
Campo Grande - MS}
\date{\today}

\begin{document}
\begin{backgroundblock}{0mm}{-10mm}
   \includegraphics[width=\paperwidth]{Figures/Modelo-Slides-FACOM.pdf}
\end{backgroundblock}

\DTMsetstyle{portuges}
\maketitle

\begin{frame}
    \begin{enumerate}
        \item Contextualização e Objetivos
        \item Metodologia
        \item Resultados
        \item Conclusão
    \end{enumerate}
\end{frame}

\section{Contextualização e Objetivos}
\begin{frame}{Contextualização e Objetivos}
    \begin{columns}
        \column{0.5\textwidth}
            \centering
            \includegraphics[scale=0.4]{Figures/res/OBI_Logo.png}
        \column{0.5\textwidth}
            \begin{itemize}
                \item Olimpíada Brasileira de Informática
                \item Modalidades: Programação e \textbf{Iniciação}
                \item Níveis 1, 2 e Júnior
            \end{itemize}
    \end{columns}
\end{frame}

\begin{frame}{}
    \begin{columns}
        \column{0.5\textwidth}
        \centering
        \includegraphics[scale=0.25]{Figures/res/Jogos_de_Lógica.png}
        \column{0.5\textwidth}
        \begin{itemize}
            \item Classifica as questões de 2003 a 2009
            \item Avaliar adequação das classes, e se necessário, propor novas
            \item Introdução do Chain-of-Thought prompt
        \end{itemize}
    \end{columns}
\end{frame}

\section{Metodologia}
\begin{frame}{Sobre as provas}
    \begin{itemize}
        \item Diferenças entre edições das olimpíadas
    \end{itemize}
    
    \begin{columns}
        \column{0.5\textwidth}
            \centering
            \includegraphics[scale=0.4]{Figures/res/OBI2003_F1N0_Q3.PNG} \\[1.5cm]
            OBI 2003
        \column{0.5\textwidth}
            \centering
            \includegraphics[scale=0.35]{Figures/res/OBI2013_F1N2_Q11.PNG}
            OBI 2013
    \end{columns}
\end{frame}

\begin{frame}{}
    \begin{columns}
        \column{0.5\textwidth}
            \centering
            \includegraphics[scale=0.4]{Figures/res/OBI2009_F2N2_Q6.PNG}
        \column{0.5\textwidth}
            \begin{itemize}
                \item \textcolor{red}{Cenário}
                \item \textcolor{green}{Regras}
                \item \textcolor{blue}{Perguntas}
            \end{itemize}
    \end{columns}
\end{frame}

\begin{frame}{Classes e Tipos}
    \begin{columns}
        \column{0.3\textwidth}
            \begin{itemize}
                \item<1-> Ordenação
                \begin{itemize}
                    \item<2-> Linear
                    \item<2-> Quadrática
                    \item<2-> Circular
                    \item<2-> Livre
                \end{itemize}
            \end{itemize}
        \column{0.3\textwidth}
            \begin{itemize}
                \item<1-> Agrupamento
                \begin{itemize}
                    \item<3-> 1-Grupo
                    \item<3-> n-Grupos
                    \item[]
                    \item[] 
                \end{itemize}
            \end{itemize}
        \column{0.4\textwidth}
            \begin{itemize}
                \item<1-> Outros
                \begin{itemize}
                    \item<4-> Grupos Ordenados
                    \item<4-> Cálculo
                    \item<4-> Definição
                \end{itemize}
            \end{itemize}
    \end{columns}
\end{frame}

\begin{frame}{Cinco princípios de prompting}
    \begin{itemize}
        \item[1.] Dar direção
        \item[2.] Especificar Formato
        \item[3.] Providenciar Exemplos
        \item[4.] Avaliar Qualidade
        \item[5.] Dividir o Trabalho
    \end{itemize}
\end{frame}

\begin{frame}{Prompts}
    \begin{itemize}
        \item[1.]<1-> Zero-shot com apenas classes
        \item[2.]<2-> Zero-shot com classes e tipos
        \item[3.]<3-> One-shot
        \item[4.]<4-> Chain-of-Thoughts
    \end{itemize}
\end{frame}

\section{Resultados}
\begin{frame}{Distribuição das questões}
    \begin{table}[h]
    \centering
    \label{tab:classes-tipos-periodos}
    \begin{tabular}{|l|cc|cc|}
    \hline
    \textbf{Tipo} & \textbf{2003–2009} & \textbf{\%} &          \textbf{2023–2024} & \textbf{\%} \\
    \hline
    Ord. – Linear       & 88    & 19,6  & 155   & 34,4  \\
    Ord. – Quadrática   & 64    & 14,3  & 19    & 4,2   \\
    Ord. – Circular     & 13    & 2,9   & 12    & 4,8   \\
    Ord. - Livre        & 38    & 8,5   & 3     & 0,7   \\
    Agrp. – 1 Grupo     & 55    & 12,2  & 8     & 1,8   \\
    Agrp. – N-Grupos    & 98    & 21,8  & 22    & 4,9   \\
    Cálculo             & 33    & 7,3   & 111   & 24,7  \\
    Grupos Ordenados    & 15    & 3,3   & 43    & 9,6   \\
    Definição           & 45    & 10,0  & 75    & 16,7  \\
    \hline
    \textbf{Total}      & 449   & 100   & 450   & 100   \\
    \hline
    \end{tabular}
    \end{table}    
\end{frame}

\begin{frame}{Resultados 2003 - 2009}
    \centering
    \begin{tikzpicture}
        \begin{axis}[
        width  = \textwidth,
        height = 6.9cm,
        major x tick style = transparent,
        ybar=2*\pgflinewidth,
        bar width=14pt,
        ymajorgrids = true,
        symbolic x coords={Classe,Classe-Tipo,One-Shot, CoT},
        xtick = data,
        scaled y ticks = false,
        enlarge x limits=0.25,
        ymin=0,
        ymax=100,
        legend style={at={(0,1)}, anchor=north west, legend cell align=left,legend columns=-1},
        nodes near coords,
        node near coords style={xshift=0.05cm,
        /utils/exec={\setbox0\hbox{\pgfmathprintnumber\pgfplotspointmeta}
        \pgfmathfloattomacro{\pgfplotspointmeta}{\F}{\M}{\E}
        \pgfmathsetmacro{\myanchor}{ifthenelse(\M*pow(10,\E-3)*2-\the\wd0>0,"east","west")}
        %\typeout{\M,\E,\the\wd0,\mysign,\myanchor}
        },
        rotate=90,anchor=\myanchor,
        /pgf/number format/.cd,fixed zerofill,precision=1},
    ]
        \addplot[style={black,fill=green,mark=none}]
            coordinates {(Classe,50.56) (Classe-Tipo,62.58) (One-Shot,67.71) (CoT,83.3)};

        \addplot[style={black,fill=red,mark=none}]
             coordinates {(Classe,49.44) (Classe-Tipo,22.72) (One-Shot,16.70) (CoT,06.01)};

        \addplot[style={black,fill=yellow,mark=none}]
             coordinates {(Classe-Tipo,14.70) (One-Shot,15.59) (CoT,10.69)};

        \legend{Acertos, Erros,Acertos parciais}
    \end{axis}
    \end{tikzpicture}
\end{frame}

\begin{frame}{Resultados 2023 - 2024}
    \centering
    \begin{tikzpicture}
        \begin{axis}[
        width  = \textwidth,
        height = 6.8cm,
        major x tick style = transparent,
        ybar=2*\pgflinewidth,
        bar width=14pt,
        ymajorgrids = true,
        symbolic x coords={Classe,Classe-Tipo,One-Shot, CoT},
        xtick = data,
        scaled y ticks = false,
        enlarge x limits=0.25,
        ymin=0,
        ymax=100,
        legend style={at={(0,1)}, anchor=north west, legend cell align=left,legend columns=-1},
        nodes near coords,
        node near coords style={xshift=0.05cm,
        /utils/exec={\setbox0\hbox{\pgfmathprintnumber\pgfplotspointmeta}
        \pgfmathfloattomacro{\pgfplotspointmeta}{\F}{\M}{\E}
        \pgfmathsetmacro{\myanchor}{ifthenelse(\M*pow(10,\E-3)*2-\the\wd0>0,"east","west")}
        %\typeout{\M,\E,\the\wd0,\mysign,\myanchor}
        },
        rotate=90,anchor=\myanchor,
        /pgf/number format/.cd,fixed zerofill,precision=1},
    ]
        \addplot[style={black,fill=green,mark=none}]
            coordinates {(Classe,47.33) (Classe-Tipo,27.11) (One-Shot,60.89) (CoT,72.89)};

        \addplot[style={black,fill=red,mark=none}]
             coordinates {(Classe,52.67) (Classe-Tipo,47.56) (One-Shot,27.56) (CoT,14.22)};

        \addplot[style={black,fill=yellow,mark=none}]
             coordinates {(Classe-Tipo,25.33) (One-Shot,11.56) (CoT,12.89)};

        \legend{Acertos, Erros,Acertos parciais}
    \end{axis}
    \end{tikzpicture}
\end{frame}

\begin{frame}{Problemas encontrados}
    \begin{columns}
        \column{0.5\textwidth}
        \centering
        \begin{itemize}
            \item Análise de imagens
        \end{itemize}
        \includegraphics[scale=0.6]{Figures/Imagens.PNG}
        \column{0.5\textwidth}
        \centering
        \begin{itemize}
            \item Agrupamento de questões
        \end{itemize}
        \includegraphics[scale=0.6]{Figures/Agrupa.PNG}
    \end{columns}
\end{frame}

\section{Conclusão}
\begin{frame}{Trabalhos futuros}
    \begin{itemize}
        \item Melhoria da precisão do prompt
        \item Modificações para permitir tratamento de imagens
        \item Análise de provas posteriores
    \end{itemize}
\end{frame}

\begin{frame}{}
    \begin{center}
        \vfill
        {\Large\bf \url{https://github.com/Sr3284/Classifica-OBI}}\\
        \vfill
    \end{center}
\end{frame}

\begin{frame}
    \begin{center}
        \vfill
        {\Large \bf Obrigado pela atenção}\\
        \vfill
        {\small Ficamos a disposição para dúvidas, críticas e sugestões}\\
        \vfill
    \end{center}
\end{frame}

\end{document}
